\chapter{\texorpdfstring{Vulnerability 5.46b: Tainted\\Length (Packet OOB Write)}{Vulnerability 5.46b: Tainted Length (Packet OOB Write)}}

\section{Executive Summary}

Vulnerability 5.46b is a \textbf{confirmed memory corruption vulnerability} where an attacker-controlled \textbf{tainted packet field (TCP Data Offset)} is used directly as a loop counter in a memory write operation, causing a \textbf{buffer overflow within the in-kernel packet buffer}. The corruption writes to adjacent \textbf{packet-local memory} (Ethernet header fields), not kernel data structures.

\begin{tcolorbox}[colback=blue!5!white,colframe=blue!75!black,title=ISO Rule Violated,fonttitle=\bfseries]
ISO-IEC Rule 5.46: ``Do not apply the sizeof operator to a polymorphic object'' (and related: tainted values in sinks)
\end{tcolorbox}

\begin{tcolorbox}[colback=red!5!white,colframe=red!75!black,title=Severity,fonttitle=\bfseries]
\textbf{MEDIUM:} Direct memory corruption confirmed, but \textbf{scope-limited to packet-buffer memory}. Overwrites adjacent packet fields (MAC addresses), not kernel structures. \textbf{No privilege escalation possible}.
\end{tcolorbox}

\subsection{Why This Vulnerability Matters}

\begin{enumerate}[label=\arabic*.]
  \item \textbf{Tainted Packet Field}: TCP Data Offset (doff) is attacker-controlled
  \item \textbf{No Validation}: Value is used directly without bounds checking
  \item \textbf{Fixed-Size Buffer}: Destination MAC address buffer is only 6 bytes
  \item \textbf{Verifier Gap}: Doesn't enforce field-level bounds checking within packet structures
  \item \textbf{Scope Limitation}: Overflow stays within packet-local memory, not kernel memory
\end{enumerate}

\section{The Vulnerable Code Pattern}

\subsection{Vulnerable Code (Simplified)}

\vspace{0.3cm}
\noindent
\fcolorbox{black!30}{gray!15}{%
  \begin{minipage}[c]{\dimexpr\textwidth-2\fboxsep-2\fboxrule\relax}
    {\small\ttfamily
    /* Fixed-size buffer for MAC destination (6 bytes standard) */\\
    unsigned char h\_dest[6];\\
    \\
    /* Get tainted length from packet TCP header */\\
    \_\_u32 tainted\_length = hdr-\textgreater{}tcp-\textgreater{}doff * 4;\\
    /* doff is 4-15 (4-bit field), so length is 20-60 bytes */\\
    bpf\_printk("[taintsink]: Tainted length \%u", tainted\_length);\\
    \\
    /* Use tainted length in a memory write loop (NO VALIDATION!) */\\
    for (i = 0; i < tainted\_length; i++) \{\\
    \phantom{for (i = 0; i < tainted\_length; i++) \{}h\_dest[i] = /* packet\_data[i] */;\\
    /* BUFFER OVERFLOW if tainted\_length > 6 */\\
    \}
    }
  \end{minipage}%
}
\vspace{0.3cm}

\subsection{Why This Is Vulnerable}

\begin{enumerate}[label=\arabic*.]
  \item \textbf{Fixed-size buffer}: \texttt{h\_dest} is \textbf{only 6 bytes} (standard Ethernet MAC address)
  \item \textbf{Tainted length}: \texttt{doff * 4} is attacker-controlled via TCP Data Offset field
  \item \textbf{Range}: TCP doff is a 4-bit field (values 4-15), so \texttt{doff * 4} = 20-60 bytes
  \item \textbf{No validation}: Code directly uses tainted value as loop counter
  \item \textbf{Overflow}: Writing 60 bytes to a 6-byte buffer = \textbf{54-byte overflow} into adjacent memory
\end{enumerate}

\section{Multi-Stage Analysis: Evidence Chain}

\subsection{Stage 1: Source Code (Patch Location)}

The vulnerability is introduced via a Git patch:
\begin{tcolorbox}[colback=gray!10, colframe=gray!10, arc=2pt, boxrule=0pt]
\small\ttfamily
\texttt{XDPs/xdp\_synproxy/patches/5\_46b\_taintsink/0001-feat-5.46*.patch}
\end{tcolorbox}

\subsection{Stage 2: Bytecode (llvm-objdump)}

Using \texttt{llvm-objdump -S xdp\_synproxy\_kern.o}, the compiled bytecode contains an \textbf{unrolled loop} writing 60 times:

\vspace{0.3cm}
\noindent
\fcolorbox{black!30}{gray!15}{%
  \begin{minipage}[c]{\dimexpr\textwidth-2\fboxsep-2\fboxrule\relax}
    {\small\ttfamily
    ;~ for (i = 0; i < tainted\_length; i++) \{\\
    ;~~ h\_dest[i] = i;\\
    56: (71) r1 = *(u8 *)(r8 +18)    \# Load loop counter i\\
    57: (73) *(u8 *)(r2 -80) = r1   \# Store at offset -80+i\\
    58: (85) call bpf\_trace\_printk  \# Signal: Wrote value at idx\\
    ;~ ... (repeated 60 times for doff=15) ...
    }
  \end{minipage}%
}
\vspace{0.3cm}

\noindent
\textbf{Interpretation}: Bytecode shows 60 store instructions (one for each value in the loop). This confirms the compiler \textbf{did not optimize away} the loop unrolling despite the unsafe length value.

\noindent
\textbf{Finding}: The full 60-iteration loop is preserved in bytecode.

\subsection{Stage 3: Verifier (bpftool prog load)}

Command:
\begin{tcolorbox}[colback=gray!10, colframe=gray!10, arc=2pt, boxrule=0pt]
\small\ttfamily
\texttt{sudo bpftool prog load xdp\_synproxy\_kern.bpf.o /sys/fs/bpf/test\_synproxy type xdp -d}
\end{tcolorbox}

\noindent
\textbf{Result}: The verifier \textbf{accepts the program without error}.

\noindent
\textbf{Finding}: Verifier does not reject the 60-byte write to a 6-byte buffer. Field-level bounds checking missing.

\subsection{Stage 4: Xlated Bytecode (bpftool prog dump xlated)}

The xlated bytecode still contains all 60 store instructions:

\vspace{0.3cm}
\noindent
\fcolorbox{black!30}{gray!15}{%
  \begin{minipage}[c]{\dimexpr\textwidth-2\fboxsep-2\fboxrule\relax}
    {\small\ttfamily
    56: (71) r1 = *(u8 *)(r8 +18)\\
    57: (73) *(u8 *)(r2 -80) = r1\\
    58: (85) call bpf\_trace\_printk\\
    ;~ ... all 60 stores preserved ...
    }
  \end{minipage}%
}
\vspace{0.3cm}

\noindent
\textbf{Finding}: Xlated bytecode shows \textbf{all 60 store instructions are preserved unchanged}. No verifier-inserted bounds checks. Vulnerability survives verification.

\subsection{Stage 5: Runtime Execution (trace\_pipe)}

Trigger the vulnerability with TCP doff=15 (maximum):

\begin{tcolorbox}[colback=gray!10, colframe=gray!10, arc=2pt, boxrule=0pt]
\small\ttfamily
python3 src/exploits/5\_46b\_taintsink/5\_46b\_overflow.py
\end{tcolorbox}

\vspace{0.3cm}
\noindent
\fcolorbox{black!30}{gray!15}{%
  \begin{minipage}[c]{\dimexpr\textwidth-2\fboxsep-2\fboxrule\relax}
    \texttt{<idle>-0 [000] ..s2. 348.737756: bpf\_trace\_printk: [taintsink]: Tainted length 60}\\
    \texttt{<idle>-0 [000] ..s2. 348.738659: bpf\_trace\_printk: [taintsink]: Wrote value at idx 0}\\
    \texttt{<idle>-0 [000] ..s2. 348.739616: bpf\_trace\_printk: [taintsink]: Wrote value at idx 1}\\
    \texttt{<idle>-0 [000] ..s2. 348.740573: bpf\_trace\_printk: [taintsink]: Wrote value at idx 2}\\
    \texttt{<idle>-0 [000] ..s2. 348.741530: bpf\_trace\_printk: [taintsink]: Wrote value at idx 3}\\
    \texttt{<idle>-0 [000] ..s2. 348.757486: bpf\_trace\_printk: [taintsink]: Wrote value at idx 59}\\
    \texttt{<idle>-0 [000] ..s2. 348.758443: bpf\_trace\_printk: [taintsink]: hdr-\textgreater{}eth-\textgreater{}h\_source differs}
  \end{minipage}%
}
\vspace{0.3cm}

\noindent
\textbf{Finding}: The trace confirms:
1. Tainted length is 60 (doff=15)
2. All 60 writes execute (indices 0-59)
3. MAC source field is corrupted (comparison detects difference)

\noindent
\textbf{Links all stages together}. Vulnerability executes at runtime.

\section{Exploitation Analysis}
\noindent
Memory corruption has been \textbf{confirmed}: the verifier accepts the program, the bytecode preserves the 60-byte loop, the xlated bytecode has no bounds guards, runtime executes all 60 writes, the attacker controls the length (via TCP data offset), and adjacent memory is corrupted with the MAC source address overwritten.

\vspace{0.2cm}
\noindent
Packet memory layout shows that the \texttt{h\_dest} buffer is \textbf{6 bytes} (indices 0-5), \texttt{h\_source} is 6 bytes (indices 6-11), and \texttt{h\_proto} is 2 bytes (indices 12-13). With \textbf{doff=15 (60-byte write)}, bytes 0-5 are written to \texttt{h\_dest} (intended), bytes 6-11 overflow to \texttt{h\_source} (first overflow), bytes 12-13 overflow to \texttt{h\_proto} (second overflow), and bytes 14-59 write beyond the packet header (third overflow).

\vspace{0.2cm}
\noindent
The impact is \textbf{scope-limited}: writes stay within the packet buffer structure (not kernel stack), the overflow cannot reach kernel memory, the attacker cannot modify task credentials or eBPF context, and while the packet may be corrupted, network processing continues with limited DoS potential.

\vspace{0.2cm}
\noindent
While \textbf{direct exploitation for kernel compromise is impossible} due to eBPF sandbox constraints---execution is isolated from kernel memory, the eBPF stack is separate from kernel stack, and corrupted data affects only packet-local memory, not kernel control flow---this vulnerability demonstrates a \textbf{critical verifier gap}: the verifier doesn't enforce field-level bounds and only checks packet-level safety, representing a \textbf{systematic gap in bounds checking} that affects all packet buffer field access patterns.

\vspace{0.2cm}
\noindent
The vulnerability is \textbf{exploitable within packet-buffer scope}, but the eBPF sandbox prevents kernel-level compromise.


\section{Artifact Evidence}

All analysis logs available in: \texttt{src/exploits/5\_46b\_taintsink/logs/}

\begin{itemize}
  \item \texttt{5\_46b\_bytecode\_analysis.txt} — Full llvm-objdump output (all 60 stores)
  \item \texttt{5\_46b\_verifier\_log.txt} — Verifier output (bpftool prog load -d)
  \item \texttt{5\_46b\_xlated\_dump.txt} — Xlated bytecode with all 60 store instructions
\end{itemize}

\noindent
PoC script: \texttt{src/exploits/5\_46b\_taintsink/5\_46b\_overflow.py}

